% ============================================================
%  REE 598: Senior Design Project I — Proposal
%  Zewail City University of Science and Technology
%  Format: Times New Roman 12pt, 1.5 spacing, chapter-based
%          section/figure/table/equation numbering
% ============================================================
% !TEX TS-program = pdflatex
\documentclass[12pt,a4paper]{report}

% ── Fonts and encoding ──────────────────────────────────────
\usepackage[T1]{fontenc}
\usepackage{mathptmx}          % Times New Roman (text + math)
\usepackage{microtype}

% ── Page geometry ───────────────────────────────────────────
\usepackage[
  top=2.5cm, bottom=2.5cm,
  left=3.0cm, right=2.5cm
]{geometry}

% ── Line spacing ────────────────────────────────────────────
\usepackage{setspace}
\onehalfspacing

% ── Mathematics ─────────────────────────────────────────────
\usepackage{amsmath,amssymb}

% ── Tables and figures ──────────────────────────────────────
\usepackage{booktabs}
\usepackage{array}
\usepackage{multirow}
\usepackage{graphicx}
\usepackage{float}
\usepackage{caption}
\captionsetup{font=small, labelfont=bf}

% ── Gantt chart ─────────────────────────────────────────────
\usepackage{pgfgantt}
\usepackage{pdflscape}
\usepackage{rotating}

% ── Lists ───────────────────────────────────────────────────
\usepackage{enumitem}

% ── Header / footer ─────────────────────────────────────────
\usepackage{fancyhdr}
\pagestyle{fancy}
\setlength{\headheight}{14pt}
\fancyhf{}
\rhead{\small REE 598 Senior Design Project I}
\lhead{\small Zewail City University of Science and Technology}
\cfoot{\thepage}
\renewcommand{\headrulewidth}{0.4pt}

% ── Hyperlinks ──────────────────────────────────────────────
\usepackage[hidelinks]{hyperref}
\usepackage{xurl}
\makeatletter
\g@addto@macro{\UrlBreaks}{\do\/\do-\do_\do.\do=\do?\do&}
\makeatother

% ── Bibliography ────────────────────────────────────────────
\usepackage{cite}

% ── Chapter-based equation/figure/table numbering ───────────
\numberwithin{equation}{chapter}
\numberwithin{figure}{chapter}
\numberwithin{table}{chapter}

% ── Section numbering depth ─────────────────────────────────
\setcounter{secnumdepth}{3}
\setcounter{tocdepth}{3}

% ── Chapter title formatting ────────────────────────────────
\usepackage{titlesec}
\titleformat{\chapter}[hang]
  {\normalfont\large\bfseries}
  {\thechapter.}{0.5em}{}
\titlespacing*{\chapter}{0pt}{12pt}{6pt}

% ── Shared project macros ───────────────────────────────────
% Shared numeric assumptions (must match project_inputs.yaml)
% -----------------------------------------------------------------------
% Benban site (retained as climate validation reference)
% -----------------------------------------------------------------------
\newcommand{\BenbanLat}{24.45}
\newcommand{\BenbanLon}{32.75}

% -----------------------------------------------------------------------
% PV module baseline (used in climate validation chapter)
% -----------------------------------------------------------------------
\newcommand{\Etamod}{0.215}
\newcommand{\WaterIntensity}{1.8}
\newcommand{\Nclean}{18}
\newcommand{\Fmargin}{0.25}

% -----------------------------------------------------------------------
% Economics
% -----------------------------------------------------------------------
\newcommand{\TariffUSD}{0.020}
\newcommand{\TariffEGP}{0.97}

% -----------------------------------------------------------------------
% TEC module parameters (TEC1-12706)
% -----------------------------------------------------------------------
\newcommand{\TECImax}{6.0}
\newcommand{\TECVmax}{15.4}
\newcommand{\TECQmax}{57.2}
\newcommand{\TECDTmax}{68}
\newcommand{\TECAlpha}{0.0534}
\newcommand{\TECR}{1.98}
\newcommand{\TECK}{0.524}
\newcommand{\TECCOPnom}{0.38}

% -----------------------------------------------------------------------
% Street-lamp scale (Scale A)
% -----------------------------------------------------------------------
\newcommand{\ScaleAPV}{80}
\newcommand{\ScaleATECn}{1}
\newcommand{\ScaleAWaterDay}{0.15}
\newcommand{\ScaleACleanArea}{0.35}
\newcommand{\ScaleABattWh}{150}

% -----------------------------------------------------------------------
% Residential scale (Scale R - villa rooftop, 5 kWp)
% -----------------------------------------------------------------------
\newcommand{\ScaleRPV}{5}
\newcommand{\ScaleRTECn}{6}
\newcommand{\ScaleRWaterDay}{1.5}
\newcommand{\ScaleRCleanArea}{28}
\newcommand{\ScaleRBattWh}{1200}

% -----------------------------------------------------------------------
% Commercial scale (Scale B - rooftop, 10 kWp)
% -----------------------------------------------------------------------
\newcommand{\ScaleBPV}{10}
\newcommand{\ScaleBTECn}{12}
\newcommand{\ScaleBWaterDay}{3.0}
\newcommand{\ScaleBCleanArea}{46}
\newcommand{\ScaleBBattWh}{2500}

% -----------------------------------------------------------------------
% Industrial scale (Scale C - ground-mount, 200 kWp, stretch)
% -----------------------------------------------------------------------
\newcommand{\ScaleCPV}{200}
\newcommand{\ScaleCPdcMW}{0.200}
\newcommand{\PdcMW}{200}
\newcommand{\ScaleCTECn}{80}
\newcommand{\ScaleCWaterDay}{40}
\newcommand{\ScaleCCleanArea}{930}
\newcommand{\ScaleCBattWh}{24000}

% -----------------------------------------------------------------------
% Climate / psychrometric operating point (Egyptian summer, conservative)
% -----------------------------------------------------------------------
\newcommand{\TambSummer}{38}
\newcommand{\RHmin}{25}
\newcommand{\RHdesign}{35}
\newcommand{\TdewDesign}{18.5}
\newcommand{\TcoldDesign}{14}
\newcommand{\ThotDesign}{58}
\newcommand{\DTdesign}{44}

% -----------------------------------------------------------------------
% Robot motion (commercial/industrial rail-guided)
% -----------------------------------------------------------------------
\newcommand{\RobotSpeed}{0.05}
\newcommand{\BrushWidth}{0.25}
\newcommand{\FlowNozzle}{0.008}


% ── Convenience unit macros (no siunitx required) ───────────
\newcommand{\degC}{$^\circ$C}
\newcommand{\Watt}{\,W}
\newcommand{\kWp}{\,kW$_\mathrm{p}$}
\newcommand{\kWh}{\,kWh}
\newcommand{\kWhL}{\,kWh/L}
\newcommand{\Lday}{\,L/day}
\newcommand{\Lh}{\,L/h}
\newcommand{\gm}{\,g/m$^3$}
\newcommand{\Wh}{\,Wh}
\newcommand{\AUnit}{\,A}
\newcommand{\VUnit}{\,V}
\newcommand{\OhmU}{\,$\Omega$}
\newcommand{\WK}{\,W/K}
\newcommand{\VK}{\,V/K}
\newcommand{\KW}{\,K/W}

% ============================================================
%  DOCUMENT
% ============================================================
\begin{document}

\fancypagestyle{plain}{%
  \fancyhf{}
  \cfoot{\thepage}
  \renewcommand{\headrulewidth}{0pt}
}

% ============================================================
%  COVER PAGE
% ============================================================
\begin{titlepage}
\thispagestyle{empty}
\begin{center}

\includegraphics[height=2.2cm]{logo.jpg}\\[0.6cm]

{\large \textbf{University of Science and Technology in Zewail City}}\\[0.1cm]
{\large Department of Renewable Energy Engineering}\\[1.2cm]

\rule{\textwidth}{1.5pt}\\[0.4cm]
{\LARGE \textbf{Autonomous Thermoelectric-Based\\[0.3em]
Atmospheric Water Harvesting Robot\\[0.3em]
for Photovoltaic Panel Cleaning}}\\[0.3cm]
{\large \textit{Multi-Scale Design for Egyptian New Cities}}\\[0.3cm]
\rule{\textwidth}{1.5pt}\\[1.0cm]

{\large \textbf{REE 598: Senior Design Project I}}\\[0.3cm]
{\large Proposal Report}\\[1.0cm]

\begin{tabular}{rl}
  \textbf{Student:}       & Rahma El Hamouly \\[0.3cm]
  \textbf{Supervisor:}    & [Supervisor Name] \\[0.3cm]
  \textbf{Date:}          & \today \\[0.3cm]
  \textbf{Academic Year:} & 2025 -- 2026 \\
\end{tabular}\\[2.0cm]

\vfill
{\small Submitted in partial fulfilment of the requirements for\\
REE 598: Senior Design Project I\\
Zewail City University of Science and Technology, Egypt}

\end{center}
\end{titlepage}

% ============================================================
%  ACKNOWLEDGMENT
% ============================================================
\chapter*{Acknowledgment}
\addcontentsline{toc}{chapter}{Acknowledgment}
\thispagestyle{plain}

The author gratefully acknowledges the guidance and continuous support of the project
supervisor throughout the conception and planning phases of this work. Appreciation is
extended to the Department of Renewable Energy Engineering at Zewail City University
of Science and Technology for providing access to laboratory facilities and computational
resources. The NASA POWER Data Access Viewer, maintained by the NASA Langley Research
Center Atmospheric Science Data Center, is acknowledged as the primary source of
meteorological data used in this study. The author thanks colleagues and laboratory
staff for constructive discussions during the preliminary design phase.

\newpage

% ============================================================
%  ABSTRACT
% ============================================================
\chapter*{Abstract}
\addcontentsline{toc}{chapter}{Abstract}
\thispagestyle{plain}

Photovoltaic (PV) panels deployed in Egyptian arid environments accumulate dust at daily
rates of 0.3--0.8\,\% per day, causing monthly power losses of 15--40\,\% without
regular cleaning. Wet cleaning --- the application of low-mineral-content water with a
physical wiping action --- is the most effective remediation method, restoring 90--98\,\%
of lost optical transmittance compared to only 70--85\,\% achieved by dry brushing alone.
However, wet cleaning at scale is constrained by the scarcity and logistical cost of
fresh water at the remote and newly urbanising Egyptian sites where solar infrastructure
is growing most rapidly.

This project designs a \textbf{fully autonomous robotic PV panel cleaning system} that
eliminates the fresh-water dependency by harvesting cleaning water directly from ambient
atmospheric moisture using \textbf{thermoelectric cooler (TEC) / Peltier modules}. The
system requires no municipal water connection and no human attendance between scheduled
maintenance intervals of weeks to months, making it directly suited to the three PV
deployment scales prevalent in Egypt's new cities: (A)~pole-mounted street-lamp panels
(\ScaleAPV\Watt), (R)~residential villa rooftops (\ScaleRPV\kWp), (B)~commercial
rooftop arrays (\ScaleBPV\kWp), and (C)~industrial ground-mount blocks
(\ScaleCPV\kWp, stretch case).

The project is structured as a two-semester Senior Design sequence. In REE~598
(Part~I), the climate feasibility of TEC-based atmospheric water harvesting is validated
against five years of NASA POWER hourly data for Benban Solar Park (Aswan, Egypt),
representing the most thermodynamically demanding Egyptian site. TEC thermodynamics are
modelled from first principles to determine condensate yield, coefficient of performance
(COP), and specific energy consumption (SEC). Three autonomous robot concept designs are
developed --- covering mechanical architecture, motion system, water treatment chain,
power supply, and control logic --- and a preliminary energy and economic analysis
quantifies the efficiency improvement and payback period. In REE~599 (Part~II), a
Scale~B commercial prototype will be fabricated and validated in REE~599; the same
architecture is downscaled for residential (Scale~R) deployment.

\newpage

% ============================================================
%  TABLE OF CONTENTS / FIGURES / TABLES
% ============================================================
\tableofcontents
\newpage
\listoffigures
\addcontentsline{toc}{chapter}{List of Figures}
\newpage
\listoftables
\addcontentsline{toc}{chapter}{List of Tables}
\newpage

% ============================================================
%  NOMENCLATURE
% ============================================================
\chapter*{Nomenclature and List of Abbreviations}
\addcontentsline{toc}{chapter}{Nomenclature and List of Abbreviations}
\thispagestyle{plain}

\section*{Roman Symbols}

\begin{tabular}{@{}lll@{}}
$A_\mathrm{pv}$        & Total PV panel area to be cleaned              & m$^2$ \\
$c_p$                  & Specific heat capacity of air                  & kJ/(kg\,K) \\
$\mathrm{COP}$         & Coefficient of performance of TEC              & --- \\
$e_s$                  & Saturation vapour pressure                     & hPa \\
$E_\mathrm{annual}$    & Annual PV energy yield                         & kWh \\
$f_\mathrm{margin}$    & Water volume safety margin                     & --- \\
$h_{fg}$               & Latent heat of vaporisation at 15\degC         & kJ/kg \\
$I$                    & TEC operating current                          & A \\
$I^*$                  & COP-maximising current                         & A \\
$K$                    & TEC thermal conductance                        & W/K \\
$\mathrm{LCOW}$        & Levelised cost of water                        & USD/L \\
$\dot{m}_\mathrm{air}$ & Air mass flow rate over cold plate             & kg/s \\
$\dot{m}_w$            & Condensate mass flow rate                      & g/h \\
$N_\mathrm{TEC}$       & Number of TEC modules                         & --- \\
$N_\mathrm{clean}$     & Cleaning events per year                       & --- \\
$P_\mathrm{in}$        & TEC electrical input power                     & W \\
$Q_c$                  & TEC cold-side heat pumping rate                & W \\
$Q_h$                  & TEC hot-side heat rejection rate               & W \\
$R$                    & TEC internal electrical resistance             & $\Omega$ \\
$R_\mathrm{th}$        & Thermal resistance                             & K/W \\
$\mathrm{RH}$          & Relative humidity                              & \% \\
$s_r$                  & Daily soiling rate                             & \%/day \\
$\mathrm{SEC}$         & Specific energy consumption                    & kWh/L \\
$T_\mathrm{amb}$       & Ambient dry-bulb temperature                   & \degC \\
$T_c$                  & TEC cold-side temperature                      & \degC \\
$T_\mathrm{dp}$        & Dew-point temperature                          & \degC \\
$T_h$                  & TEC hot-side temperature                       & \degC \\
$T_\mathrm{int}$       & Cleaning interval                              & days \\
$T_m$                  & Mean TEC temperature $(T_h+T_c)/2$             & K \\
$V_\mathrm{event}$     & Water volume per cleaning event                & L \\
$\dot{V}_w$            & Condensate volumetric yield                    & L/h \\
$w$                    & Water application intensity                    & L/m$^2$ \\
$Z$                    & TEC figure of merit $\alpha^2/(RK)$            & K$^{-1}$ \\
\end{tabular}

\vspace{1em}
\section*{Greek Symbols}

\begin{tabular}{@{}lll@{}}
$\alpha$               & Seebeck coefficient                            & V/K \\
$\Delta T$             & Temperature differential $T_h - T_c$           & K \\
$\eta_c$               & Condensate collection efficiency               & --- \\
$\eta_\mathrm{mod}$    & PV module conversion efficiency                & --- \\
$\eta_\mathrm{PR}$     & PV system performance ratio                    & --- \\
$\rho_v$               & Absolute humidity                              & g/m$^3$ \\
$\theta_\mathrm{sa}$   & Heat-sink-to-ambient thermal resistance        & K/W \\
$\pi$                  & Electricity tariff                             & USD/kWh \\
\end{tabular}

\vspace{1em}
\section*{Abbreviations}

\begin{tabular}{@{}ll@{}}
AWG    & Atmospheric Water Generator \\
AWH    & Atmospheric Water Harvesting \\
BLDC   & Brushless Direct-Current (motor) \\
BOM    & Bill of Materials \\
CAE    & Computer-Aided Engineering \\
COP    & Coefficient of Performance \\
DoD    & Depth of Discharge \\
GHI    & Global Horizontal Irradiance \\
HDPE   & High-Density Polyethylene \\
IP     & Ingress Protection \\
LCOW   & Levelised Cost of Water \\
LiFePO$_4$ & Lithium Iron Phosphate (battery chemistry) \\
MPPT   & Maximum Power Point Tracking \\
NAC    & New Administrative Capital (Egypt) \\
PM     & Particulate Matter \\
PV     & Photovoltaic \\
RH     & Relative Humidity \\
RTC    & Real-Time Clock \\
SEC    & Specific Energy Consumption \\
SoC    & State of Charge \\
TDS    & Total Dissolved Solids \\
TEC    & Thermoelectric Cooler (Peltier module) \\
UV     & Ultraviolet \\
VCR    & Vapour-Compression Refrigeration \\
\end{tabular}
\newpage

% ============================================================
%  CHAPTER 1 — INTRODUCTION
% ============================================================
\chapter{Introduction}

\section{Importance of the Project Field}

Photovoltaic energy is central to Egypt's long-term energy strategy. Egypt's Vision 2030
Sustainable Development Strategy targets 42\,\% of electricity from renewable sources by
2035, with solar PV constituting the largest share \cite{egypt2030vision}. With an annual
global horizontal irradiance (GHI) of 5.0--6.5\,kWh/m$^2$/day --- among the highest in
the world --- Egypt has both the resource and the policy mandate to aggressively expand PV
capacity \cite{globalsolaratlas,irena2023egypt}. Installed utility-scale PV capacity
reached 4.1\,GW by end-2023 and a further 6\,GW pipeline is under active development
\cite{irena2023egypt}.

Egypt is simultaneously building a network of entirely new cities. The New Administrative
Capital (NAC), New Alamein, New Mansoura, and more than a dozen others integrate distributed
solar energy at the urban-planning level: PV-powered street lighting, commercial rooftop
installations, and large ground-mount plants serving the municipal grid
\cite{egypt2030vision}. These applications span three orders of magnitude in installed
capacity --- from individual 80\,W street-lamp panels to 200\,kWp industrial blocks ---
and share the same fundamental operational challenge: the progressive degradation of panel
output by dust soiling.

\section{Definition of Main Challenges in the Field}

\subsection{Dust Soiling and Energy Loss}

In Egyptian arid environments, desert dust, PM\textsubscript{10} particulate matter,
calcium carbonate aerosols, and biological deposits accumulate on the glass surface of PV
modules. These deposits scatter and absorb incident solar radiation, reducing optical
transmittance and degrading electrical output. Field measurements across Upper Egypt
and equivalent MENA arid sites consistently report daily soiling rates of
0.3--0.8\,\%/day under normal conditions, with monthly accumulated losses of 15--40\,\%
during the dry season (April--September) \cite{kazem2023soiling,sarver2013pvsoiling}.
Episodic khamsin dust-storm events can deposit up to 10\,\% loss in a single day
\cite{caron2013soilingbenban}.

\subsection{Limitations of Existing Cleaning Methods}

Manual wet cleaning is effective but operationally intensive and expensive at scale.
Automated dry-brush robots reduce labour but achieve only 70--85\,\% transmittance
restoration \cite{mejia2021cleaning,costa2021pvrobot}. Pressurised water skids are
effective but require a continuous water supply unavailable at remote or newly developing
sites. No commercially available autonomous PV cleaning system integrates its own
independent water source.

\subsection{Water Scarcity at New-City Deployment Sites}

Egypt's new cities are built in locations where municipal water infrastructure matures
years after solar panel installation. Street-lamp PV units are distributed across urban
footprints of tens of square kilometres; commercial arrays sit on buildings under
construction; industrial plants are at desert peripheries. Supplying cleaning water by
tanker increases operational cost, carbon footprint, and logistical complexity --- all
contrary to the sustainability mandate of these cities.

\section{Problem or Challenge to Be Solved}

\textit{How to clean PV panels continuously and autonomously at Egyptian new-city
deployment sites --- across street-lamp, residential, commercial, and (where
applicable) industrial scales --- without
relying on a piped or delivered water supply, and without requiring regular human
attendance?}

The proposed solution is an \textbf{autonomous PV cleaning robot that harvests its own
cleaning water from atmospheric moisture} using thermoelectric cooler (TEC) modules,
powered entirely from the PV array it maintains.

% ============================================================
%  CHAPTER 2 — BACKGROUND AND LITERATURE REVIEW
% ============================================================
\chapter{Background and Literature Review}

\section{Why Wet Cleaning Outperforms Dry Methods}

\subsection{Dust Adhesion and the Role of Water}

Dust particles adhere to module glass through van der Waals forces, electrostatic
interaction, and --- for coarse particles --- gravitational interlocking with surface
micro-irregularities \cite{sarver2013pvsoiling}. In Egyptian environments, calcium
carbonate particles form a cementing crust under overnight humidity cycling, significantly
increasing adhesion force \cite{kazem2023soiling}. Dry brushing dislodges weakly-bonded
particles but is ineffective against cemented or electrostatically-bound fine particles
($<$10\,\textmu m). Water dissolves the carbonate cement, reduces electrostatic forces
through ionic screening, and enables removal of all particle size classes
\cite{mejia2021cleaning}.

\subsection{Transmittance Restoration Comparison}

Table~\ref{tab:cleaning_comparison} summarises transmittance restoration efficiencies
reported across the literature for different cleaning methods.

\begin{table}[H]
\centering
\caption{Transmittance restoration by cleaning method (literature summary).}
\label{tab:cleaning_comparison}
\begin{tabular}{lcc}
\toprule
\textbf{Cleaning Method} & \textbf{Restoration (\%)} & \textbf{Reference} \\
\midrule
Dry brush / wiper only               & 70--85 & \cite{mejia2021cleaning,costa2021pvrobot} \\
Dry pressurised air blowing          & 60--78 & \cite{kazem2023soiling} \\
Wet spray + manual wipe              & 90--98 & \cite{mejia2021cleaning} \\
Wet automated brush + deionised water & 93--99 & \cite{beattie2012waterspots} \\
No cleaning (soiling accumulation)   & 0      & --- \\
\bottomrule
\end{tabular}
\end{table}

\noindent The 15--28 percentage-point advantage of wet over dry cleaning translates
directly to recoverable energy: for a 10\,kWp array with a 15\,\% annual soiling loss,
shifting from dry to wet cleaning recovers an additional 2--4\,\% of annual energy output.

\subsection{Importance of Low-TDS Rinse Water and the TEC Advantage}

Mineral-rich water (TDS $>$200\,mg/L) leaves a residual film of calcium and magnesium
carbonate on the glass surface after evaporation --- the ``water-spot'' effect. Beattie
et al.\ \cite{beattie2012waterspots} demonstrated that water spots from high-TDS tap
water can reduce transmittance by 3--7\,\%, partially negating the cleaning benefit.
\textbf{TEC-condensed water is distilled quality (TDS $<$ 5\,mg/L)}, eliminating the
water-spot problem and enabling spot-free, high-efficiency cleaning. This is a decisive
advantage of the atmospheric source over tap or tanker-supplied water.

\section{Atmospheric Water Harvesting Technologies}

\subsection{Cooling Condensation: TEC vs.\ VCR}

Active condensation systems cool a surface below the local dew-point temperature.
Vapour-compression refrigeration (VCR) achieves a specific energy consumption (SEC) of
0.3--0.8\,kWh/L at RH above 50\,\% \cite{bagheri2024sec}. Thermoelectric cooler (TEC)
systems achieve 2--12\,kWh/L depending on operating conditions
\cite{alqaed2023awg,eslami2021tec}. While TEC has higher SEC than VCR, its compact form
factor, absence of refrigerant, direct DC compatibility, and negligible maintenance
requirements make it the preferred technology for an embedded, mobile robot application.
The robot draws energy from the PV array at zero marginal cost, so SEC is not a binding
economic constraint.

\subsection{Desiccant Sorption and MOF Systems}

Solid desiccants (silica gel, zeolite, metal-organic frameworks) adsorb atmospheric
moisture passively and release it thermally \cite{zhou2020review,feng2022review}. The
cyclic nature, large sorption-bed mass, and complex thermal management make integration
into a mobile cleaning robot impractical for the current application.

\subsection{Passive Dew and Fog Collection}

At the low night-time RH values characteristic of inland Egyptian sites (20--35\,\%),
passive dew/fog yields are insufficient to meet the cleaning water demand
\cite{feng2022review}.

\section{TEC Module Performance in Arid Climates: Literature Evidence}

Shourideh et al.\ \cite{shourideh2021tec} characterised a TEC1-12706-based AWG at
$T_\mathrm{amb}$ = 25--40\,\degC\ and RH = 30--65\,\%, achieving condensate rates of
0.06--0.20\,L/h per module. Eslami et al.\ \cite{eslami2021tec} validated a thermal
model of a Peltier AWG and established that the cold-side temperature must be at least
3--5\,K below the dew-point for reliable condensation. Alqaed et al.\
\cite{alqaed2023awg} demonstrated a solar-powered TEC AWG in Saudi Arabian arid
conditions ($T_\mathrm{amb} = 35$\,\degC, RH = 30--55\,\%) achieving 0.10--0.50\,L/h
per module and SEC of 2.8--9.4\,kWh/L --- conditions comparable to Benban.

\section{Autonomous PV Cleaning Robots: State of the Art}

Costa et al.\ \cite{costa2021pvrobot} catalogue both commercial and research-prototype
robotic cleaning systems. All commercial robots (e.g., Ecoppia, Heliotex, Miraikikai)
require an external water connection; none integrates on-board water harvesting. Hasan
et al.\ \cite{hasan2023cleaningrobot} highlight the operational gap of water supply at
remote sites. Pan et al.\ \cite{pan2022robotdesign} determine optimum cleaning intervals
by balancing soiling loss against cleaning cost; their methodology is adapted here for
Egyptian conditions.

\section{Critique of Prior Approaches and the Research Gap}

Prior AWH-for-PV-cleaning work correctly identified TEC feasibility but stopped at the
water-supply calculation without designing the cleaning robot. Robotic cleaning studies
assume water is available and do not address supply. \textbf{This project is the first to
integrate the complete chain --- from atmospheric moisture to clean panel glass --- into a
single autonomous, self-powered design}. The use of distilled-quality TEC condensate
also eliminates the water-spot problem inherent in prior wet-cleaning implementations
that use mineralised water sources.

% ============================================================
%  CHAPTER 3 — PROBLEM DEFINITION
% ============================================================
\chapter{Problem Definition}

\section{Problem Statement}

PV panels in Egyptian arid environments lose 15--40\,\% of monthly energy output to dust
soiling. Wet cleaning is the most effective remediation but requires fresh, low-TDS water.
At the four dominant PV deployment scales in Egypt's new cities --- street-lamp,
residential villa rooftop, commercial rooftop, and industrial ground-mount ---
piped or delivered water is either
unavailable, cost-prohibitive, or logistically impractical. Current autonomous cleaning
robots carry no independent water source, leaving the cleaning duty unmet without human
attendance and an external water supply chain.

\section{Proposed Solution Architecture}

The proposed solution is a \textbf{self-contained autonomous robot system} comprising
five integrated subsystems:

\begin{enumerate}[leftmargin=*, label=(\roman*)]
  \item A \textbf{TEC-based atmospheric water harvester} that condenses moisture from
        ambient night-time air onto a Peltier-cooled plate. No water supply connection
        required.
  \item A \textbf{compact water treatment train} (sediment filter $\rightarrow$ activated
        carbon $\rightarrow$ UV-C sterilisation) producing rinse-quality water
        (TDS $<$ 50\,mg/L from the starting TDS $<$ 5\,mg/L condensate).
  \item A \textbf{mechanical cleaning robot} that traverses the panel surface, delivering
        water through spray nozzles ahead of a microfibre brush/squeegee bar, restoring
        90--98\,\% of optical transmittance per cleaning event.
  \item A \textbf{PV self-power system} (MPPT charge controller + LiFePO\textsubscript{4}
        battery) drawing energy from the very array being maintained, consuming only
        2--10\,\% of daily output.
  \item An \textbf{embedded control system} (ESP32-class microcontroller) executing
        autonomous harvesting, scheduling, traversal, and safety logic with no human
        input required between maintenance intervals of 1--3 months.
\end{enumerate}

\section{Deployment Scales and Main Specifications}

Four deployment scales are defined, in order of deployment priority for Egyptian
new cities: street-lamp (A), residential villa rooftop (R), commercial rooftop (B),
and industrial ground-mount (C, stretch case). Design parameters are summarised in
Table~\ref{tab:specs}.

\begin{table}[H]
\centering
\caption{Design specifications by deployment scale (priority order: A\,$\to$\,R\,$\to$\,B\,$\to$\,C).}
\label{tab:specs}
\small
\begin{tabular}{p{4.0cm}cccc}
\toprule
\textbf{Parameter} & \textbf{Scale A} & \textbf{Scale R} & \textbf{Scale B} & \textbf{Scale C} \\
                   & \textit{Street lamp} & \textit{Residential} & \textit{Commercial} & \textit{Industrial$^\dagger$} \\
\cmidrule(lr){2-5}
PV system rating                   & \ScaleAPV\,W  & \ScaleRPV\,kWp  & \ScaleBPV\,kWp  & \ScaleCPV\,kWp \\
Panel area (m$^2$)                 & \ScaleACleanArea  & \ScaleRCleanArea  & \ScaleBCleanArea  & \ScaleCCleanArea \\
Water per cleaning event (L)       & 0.63   & 50    & 82.8   & 1674 \\
Target condensate per day (L)      & \ScaleAWaterDay & \ScaleRWaterDay & \ScaleBWaterDay & \ScaleCWaterDay \\
TEC modules required               & \ScaleATECn & \ScaleRTECn & \ScaleBTECn & \ScaleCTECn \\
Battery capacity (Wh)              & \ScaleABattWh & \ScaleRBattWh & \ScaleBBattWh & \ScaleCBattWh \\
Cleaning frequency                 & Every 14 days & Every 14 days & Every 14 days & Every 7--14 days \\
Robot mass limit (kg)              & $\leq$2  & $\leq$8  & $\leq$15  & $\leq$60 \\
Motion system                      & Pole-rail wiper & Roof-rail rover & Rail-guided carriage & Gantry rail vehicle \\
Human maintenance interval         & 3 months  & 2 months  & 1 month   & 2 weeks \\
\bottomrule
\end{tabular}
\vspace{0.25em}
\par\noindent\footnotesize $\dagger$Scale~C retained as an upwards-scalable stretch case; the primary design focus is Scales A/R/B.\normalsize
\end{table}

\noindent Water per cleaning event is derived from the panel area and design water
intensity:
\begin{equation}
  V_\mathrm{event} = w \cdot A_\mathrm{pv}, \qquad w = \WaterIntensity\;\mathrm{L/m^2}.
  \label{eq:3.1}
\end{equation}
The daily water production target, accounting for the cleaning interval $T_\mathrm{int}$
and safety margin $f_\mathrm{margin}$:
\begin{equation}
  \dot{V}_\mathrm{day} = \frac{V_\mathrm{event}}{T_\mathrm{int}}(1 + f_\mathrm{margin}),
  \qquad T_\mathrm{int} = 14\;\mathrm{days},\quad f_\mathrm{margin} = \Fmargin.
  \label{eq:3.2}
\end{equation}

% ============================================================
%  CHAPTER 4 — OBJECTIVES AND OUTCOMES
% ============================================================
\chapter{Objectives and Outcomes}

\section{Objectives of REE 598 (Senior Design Project I)}

\begin{enumerate}[leftmargin=*, label=\textbf{O\arabic*.}]
  \item \textbf{Climate feasibility validation.} Demonstrate through quantitative
        psychrometric analysis of five-year NASA POWER hourly data at Benban Solar Park
        that TEC modules can achieve the required condensate yields under the most adverse
        Egyptian arid conditions.
  \item \textbf{TEC thermodynamic sizing.} Derive from first principles the COP, cold-side
        temperature, condensate yield, and SEC of TEC modules at the design operating
        point; validate against published experimental benchmarks.
  \item \textbf{System architecture design.} Develop complete autonomous robot system
        designs for all four deployment scales (A/R/B/C), covering mechanical, water,
        power, and control subsystems, with primary focus on the street-lamp,
        residential, and commercial scales.
  \item \textbf{Wet-cleaning efficiency quantification.} Calculate the improvement in
        annual PV energy yield achievable by the proposed wet-cleaning regime relative to
        baseline (no cleaning) and dry-only cleaning, using a validated soiling loss model.
  \item \textbf{Preliminary economic analysis.} Estimate capital cost, annual energy
        recovered, and simple payback period for each scale.
  \item \textbf{Prototype preparation.} Produce a complete bill of materials, fabrication
        plan, and validation protocol for the Scale~B prototype, ready for construction
        in REE~599.
\end{enumerate}

\section{Expected Outcomes of REE 598}

\begin{itemize}[leftmargin=*]
  \item Psychrometric feasibility analysis with monthly condensation-window data
        (tabulated and plotted from Python simulation).
  \item TEC sizing calculations with all intermediate steps, validated against
        \cite{eslami2021tec,alqaed2023awg,shourideh2021tec}.
  \item Four-scale robot concept designs (A/R/B/C) with dimensioned sketches and design
        justification.
  \item Quantified annual energy recovery table comparing baseline, dry-clean, and
        wet-clean scenarios for each scale.
  \item Economic summary table: capital cost, annual benefit, simple payback.
  \item Scale~B bill of materials (BOM), fabrication timeline, and validation test plan.
  \item This proposal report and a technical presentation.
\end{itemize}

% ============================================================
%  CHAPTER 5 — ASSUMPTIONS AND CONSTRAINTS
% ============================================================
\chapter{Assumptions and Constraints}

\section{Design Assumptions}

\begin{enumerate}[leftmargin=*, label=\textbf{A\arabic*.}]
  \item Benban Solar Park (\BenbanLat$^\circ$N, \BenbanLon$^\circ$E) is the climate
        validation reference site. All new-city deployment locations have equal or more
        favourable humidity conditions.
  \item Conservative summer design point: $T_\mathrm{amb}$ = \TambSummer\,\degC,
        RH = \RHdesign\,\%, night harvesting window = 6\,h.
  \item TEC reference module: TEC1-12706 ($I_\mathrm{max}$ = \TECImax\,A,
        $V_\mathrm{max}$ = \TECVmax\,V, $Q_\mathrm{max}$ = \TECQmax\,W at
        $T_h$ = 27\,\degC, $\Delta T_\mathrm{max}$ = \TECDTmax\,K)
        \cite{eslami2021tec}.
  \item Condensate collection efficiency $\eta_c$ = 0.80 (20\,\% drainage and
        evaporation losses).
  \item Soiling loss model: linear accumulation at 0.5\,\%/day (median for Upper Egypt
        \cite{kazem2023soiling}), reset to zero after each wet cleaning event.
  \item Water application intensity $w$ = \WaterIntensity\,L/m$^2$ per cleaning event
        \cite{mejia2021cleaning}.
  \item PV module efficiency $\eta_\mathrm{mod}$ = \Etamod\ (monocrystalline Si,
        standard commercial grade).
  \item Battery chemistry: LiFePO\textsubscript{4}, depth of discharge (DoD)
        limited to 80\,\%.
  \item Robot is self-powered exclusively from the maintained PV array via MPPT
        charge controller; no grid connection assumed.
\end{enumerate}

\section{Constraints}

\begin{enumerate}[leftmargin=*, label=\textbf{C\arabic*.}]
  \item \textbf{Contact force:} brush normal force on panel glass $<$ 5\,N
        (IEC 61215 structural handling limit).
  \item \textbf{Water quality:} post-treatment TDS $<$ 50\,mg/L to prevent
        water-spot residue \cite{beattie2012waterspots}.
  \item \textbf{IP rating:} all electronics and TEC enclosures rated IP54 minimum
        per IEC 60529.
  \item \textbf{Robot mass:} not to exceed the panel frame point-load rating
        ($\approx$200\,N/m for standard aluminium frames).
  \item \textbf{Energy self-sufficiency:} robot 24-hour energy demand $\leq$10\,\% of
        daily PV array output.
\end{enumerate}

\section{Relevant Regulations, Codes, and Standards}

\begin{itemize}[leftmargin=*]
  \item \textbf{IEC 61215:2021} --- Crystalline silicon PV modules; structural load and
        mechanical stress requirements.
  \item \textbf{IEC 62446-1:2016} --- PV systems; testing, documentation and maintenance.
  \item \textbf{IEC 60529:2013} --- Degrees of protection (IP Code).
  \item \textbf{EgyptERA Tariff Circular, August 2024} --- Electricity tariff applicable
        to self-consumption \cite{egyptera_tariff}.
  \item \textbf{Egyptian National Building Code, Part 7} --- Renewable energy system
        installation requirements for new cities.
\end{itemize}

\section{Technical, Economic, Environmental, and Ethical Considerations}

\subsection{Technical}

TEC modules have a thermal-cycling life of 50,000--200,000 cycles depending on $\Delta T$.
Operating at 70\,\% of $I_\mathrm{max}$ reduces $\Delta T$ and extends module life to a
10-year design horizon without replacement.

\subsection{Economic}

Capital cost is dominated by the LiFePO\textsubscript{4} battery ($\approx$30\,\% of BOM).
Battery prices fell approximately 40\,\% between 2022 and 2024 and continue to decline,
improving the economic case \cite{irena2023egypt}. The lack of water supply cost further
reduces operational expenditure relative to all competing systems.

\subsection{Environmental}

The robot is fully self-powered from renewable energy, uses no municipal water, and
generates no chemical waste. TEC modules contain bismuth telluride
(Bi\textsubscript{2}Te\textsubscript{3}), a mildly toxic heavy-metal compound;
end-of-life disposal must comply with Egyptian e-waste regulations.

\subsection{Ethical and Societal}

Autonomous operation reduces dependence on manual cleaning labour. The project
team commits to assessing the employment impact on cleaning contractors in new-city
PV projects and recommending a managed transition pathway.

% ============================================================
%  CHAPTER 6 — BODY: MATHEMATICAL MODELLING AND METHODOLOGY
% ============================================================
\chapter{Mathematical Modelling and Methodology}

\section{Soiling Loss Model and Wet-Cleaning Efficiency}

\subsection{Soiling Accumulation and Reset}

The soiling loss factor $L_s$ evolves linearly between cleaning events:
\begin{equation}
  L_s(t_d) = s_r \cdot t_d, \qquad 0 \leq t_d \leq T_\mathrm{int},
  \label{eq:6.1}
\end{equation}
where $s_r = 0.5$\,\%/day and $t_d$ is days since the last cleaning. At the moment of
wet cleaning, $L_s$ resets to zero (99\,\% transmittance restoration efficiency assumed
for TEC distilled-water cleaning). The time-average soiling loss over an interval is:
\begin{equation}
  \bar{L}_s = \frac{s_r \cdot T_\mathrm{int}}{2} = \frac{0.5 \times 14}{2} = 3.5\;\%.
  \label{eq:6.2}
\end{equation}

\subsection{Annual Energy Recovery}

For dry-only cleaning (70--85\,\% restoration), a residual post-clean soiling of
0.7\,\% per interval remains. The additional annual energy recovered by switching from
dry to wet cleaning:
\begin{equation}
  \Delta E_\mathrm{wet-dry} = E_\mathrm{annual} \cdot \eta_\mathrm{PR} \cdot
  \bigl[\bar{L}_{s,\mathrm{dry\,residual}} - \bar{L}_{s,\mathrm{wet\,residual}}\bigr],
  \label{eq:6.3}
\end{equation}
where $\bar{L}_{s,\mathrm{wet\,residual}} \approx 0$ (distilled water, no spot residue)
and $\bar{L}_{s,\mathrm{dry\,residual}} \approx 0.70\,\%$.

Table~\ref{tab:energy_recovery} quantifies the annual energy and financial benefit for
each deployment scale.

\begin{table}[H]
\centering
\caption{Annual energy recovery: wet cleaning vs.\ dry-only and no cleaning.}
\label{tab:energy_recovery}
\begin{tabular}{lccccc}
\toprule
\textbf{Scale} & \textbf{Rating} & \textbf{Annual yield} &
\multicolumn{2}{c}{\textbf{Energy recovered by wet clean}} & \textbf{Financial value} \\
\cmidrule(lr){4-5}
& & (kWh) & vs.\ no clean (kWh) & vs.\ dry (kWh) & vs.\ dry (USD/yr) \\
\midrule
A & 80\,W     & 135       & 27.0      & 0.95    & 0.02 \\
R & 5\,kWp    & 8{,}500   & 1{,}700   & 59      & 1.18 \\
B & 10\,kWp   & 16{,}800  & 3{,}360   & 118     & 2.36 \\
C & 200\,kWp  & 336{,}000 & 67{,}200  & 2{,}352 & 47.0 \\
\bottomrule
\end{tabular}
\vspace{0.25em}
\par\noindent\footnotesize Tariff: \TariffUSD\,USD/kWh (EgyptERA industrial, Aug.\ 2024 \cite{egyptera_tariff}). Annual yield at 5.8\,kWh/m$^2$/day, PR = 0.80 \cite{globalsolaratlas}. ``vs.\ no clean'' assumes a 20\,\% untreated annual soiling loss; ``vs.\ dry'' isolates the incremental gain of wet over dry cleaning (0.7 percentage-point residual).\normalsize
\end{table}

\noindent The energy recovered versus no cleaning (20\,\% average annual soiling loss
assumed) is substantially larger; the table isolates the \textit{incremental} gain from
choosing wet over dry to highlight the specific contribution of the TEC water supply.

\section{TEC Thermodynamic Model}

\subsection{Governing Peltier--Seebeck--Fourier Equations}

A single-stage TEC module is governed by \cite{riffat2003peltier,zhao2022teccondenser}:
\begin{align}
  Q_c       &= \alpha T_c I - \tfrac{1}{2}I^2 R - K(T_h - T_c), \label{eq:6.4} \\
  P_\mathrm{in} &= \alpha(T_h - T_c)I + I^2 R,                  \label{eq:6.5} \\
  Q_h       &= Q_c + P_\mathrm{in},                              \label{eq:6.6} \\
  \mathrm{COP} &= Q_c \,/\, P_\mathrm{in}.                       \label{eq:6.7}
\end{align}
Module parameters for TEC1-12706 \cite{eslami2021tec}:
$\alpha = \TECAlpha$\,V/K, $R = \TECR$\,$\Omega$, $K = \TECK$\,W/K.

\subsection{Dew-Point Requirement}

For reliable condensation, the cold-side temperature must satisfy \cite{eslami2021tec}:
\begin{equation}
  T_c \leq T_\mathrm{dp} - 4\;\mathrm{K}.
  \label{eq:6.8}
\end{equation}
The dew point is computed via the Magnus approximation \cite{bagheri2024sec}:
\begin{equation}
  T_\mathrm{dp} = \frac{243.04\left[\ln\!\left(\frac{\mathrm{RH}}{100}\right)
    + \dfrac{17.625\,T_\mathrm{amb}}{243.04 + T_\mathrm{amb}}\right]}
    {17.625 - \ln\!\left(\frac{\mathrm{RH}}{100}\right)
    - \dfrac{17.625\,T_\mathrm{amb}}{243.04 + T_\mathrm{amb}}}.
  \label{eq:6.9}
\end{equation}
At design conditions ($T_\mathrm{amb}$ = \TambSummer\,\degC, RH = \RHdesign\,\%):
$T_\mathrm{dp}$ = \TdewDesign\,\degC, so $T_c \leq$ \TcoldDesign\,\degC,
requiring $\Delta T = T_h - T_c \leq$ \DTdesign\,K ---
well within the module's rated $\Delta T_\mathrm{max}$ = \TECDTmax\,K.

\subsection{COP-Optimising Current}

The current maximising COP is \cite{luo2022tecmodel}:
\begin{equation}
  I^* = \frac{\sqrt{1 + ZT_m} - 1}{\alpha/K},
  \label{eq:6.10}
\end{equation}
where $Z = \alpha^2/(RK)$ is the figure of merit and $T_m = (T_h + T_c)/2$.
\begin{gather*}
  Z = \frac{(\TECAlpha)^2}{\TECR \times \TECK}
    = 2.76 \times 10^{-3}\;\mathrm{K^{-1}}, \qquad
  T_m = \frac{\ThotDesign + \TcoldDesign}{2} + 273.15 = 309\;\mathrm{K}.
\end{gather*}
Evaluating Eq.~\eqref{eq:6.10}: $I^* \approx$ 3.2\,A. Substituting into
Eqs.~\eqref{eq:6.4}--\eqref{eq:6.7}:
\begin{equation*}
  Q_c \approx 8.4\;\mathrm{W}, \quad
  P_\mathrm{in} \approx 22.1\;\mathrm{W}, \quad
  \mathrm{COP} \approx \TECCOPnom.
\end{equation*}

\subsection{Condensate Yield per Module}

The condensate yield follows from the humid-air mass balance across the cold plate
\cite{eslami2021tec,shourideh2021tec}:
\begin{equation}
  \dot{V}_w = \eta_c \cdot \dot{m}_\mathrm{air} \cdot
  (\rho_{v,\mathrm{in}} - \rho_{v,\mathrm{out}}) \times 10^{-3} \quad\mathrm{[L/h]},
  \label{eq:6.11}
\end{equation}
where absolute humidity is computed from:
\begin{equation}
  \rho_v = \frac{0.622 \cdot e_s(T) \cdot \mathrm{RH}/100}
           {0.2871\,(T + 273.15)} \times 10^3 \quad\mathrm{[g/m^3]},
  \quad e_s(T) = 6.112\exp\!\left(\frac{17.625\,T}{243.04+T}\right).
  \label{eq:6.12}
\end{equation}
At design conditions: $\rho_{v,\mathrm{in}}$ = 14.7\,g/m$^3$,
$\rho_{v,\mathrm{out}}$ = 11.4\,g/m$^3$ (saturated at $T_c$ = 14\,\degC).
With $\dot{m}_\mathrm{air}$ = 2.71$\times$10$^{-3}$\,kg/s and $\eta_c$ = 0.80:
\begin{equation*}
  \dot{V}_w \approx 0.026\;\mathrm{L/h\,per\,module}.
\end{equation*}
Over a 6-hour summer night: $\approx$0.155\,L/module/day (satisfies Scale~A target
directly; Scales~B and C use multiple modules as per Table~\ref{tab:specs}).

\subsection{Specific Energy Consumption}

\begin{equation}
  \mathrm{SEC} = \frac{P_\mathrm{in} + P_\mathrm{fan}}{\dot{V}_w}
  = \frac{22.1 + 2.0}{0.026} \approx 0.927\;\mathrm{kWh/L}.
  \label{eq:6.13}
\end{equation}

\section{Heat Sink Sizing}

The hot-side heat rejection $Q_h \approx$ 30.5\,W must be dissipated to ambient.
Required thermal resistance:
\begin{equation}
  \theta_\mathrm{sa} = \frac{T_h - T_\mathrm{amb}}{Q_h}
  = \frac{\ThotDesign - \TambSummer}{30.5}
  \approx 0.66\;\mathrm{K/W}.
  \label{eq:6.14}
\end{equation}
Standard 50$\times$50\,mm extruded aluminium heat sinks with a 60\,mm fan achieve
$\theta_\mathrm{sa} \approx$ 0.5\,K/W, satisfying Eq.~\eqref{eq:6.14}
\cite{lin2022heatsink}.

\section{Power Budget and Battery Sizing (Scale B)}

\begin{align}
  E_\mathrm{TEC}   &= N_\mathrm{TEC} \cdot P_\mathrm{in} \cdot t_\mathrm{harvest}
    = 12 \times 24.1\;\mathrm{W} \times 6\;\mathrm{h} = 1735\;\mathrm{Wh}, \label{eq:6.15}\\
  E_\mathrm{robot} &= (P_\mathrm{drive} + P_\mathrm{pump}) \cdot t_\mathrm{clean}
    = 160\;\mathrm{W} \times 0.5\;\mathrm{h} = 80\;\mathrm{Wh},            \label{eq:6.16}\\
  E_\mathrm{ctrl}  &= 5\;\mathrm{W} \times 24\;\mathrm{h} = 120\;\mathrm{Wh},      \label{eq:6.17}\\
  E_\mathrm{total} &\approx 1935\;\mathrm{Wh}.                                      \label{eq:6.18}
\end{align}
Battery capacity at DoD = 80\,\%:
\begin{equation}
  C_\mathrm{batt} = \frac{E_\mathrm{total}}{0.80}
  = \frac{1935}{0.80} \approx 2420\;\mathrm{Wh}
  \;\rightarrow\; \ScaleBBattWh\;\mathrm{Wh\;(nominal,\;with\;7\,\%\;margin)}.
  \label{eq:6.19}
\end{equation}
Self-consumption: $1935\,\mathrm{Wh} / 58{,}000\,\mathrm{Wh} = 3.3\,\%$
--- well within constraint C5. The same sizing procedure, applied with $N_\mathrm{TEC}=6$
and $E_\mathrm{TEC}\approx 870$\,Wh, yields a Scale~R battery of \ScaleRBattWh\,Wh;
street-lamp Scale~A requires only \ScaleABattWh\,Wh with a single module.

\section{Proposed Methodology}

\subsection{Phase 1 --- Climate and TEC Feasibility Validation}

Five-year (2019--2023) NASA POWER hourly data \cite{nasa_power} for Benban are processed
by \texttt{scripts/psychrometrics\_hourly.py} to compute hourly $T_\mathrm{dp}$ and
$\rho_v$ via Eqs.~\eqref{eq:6.9} and \eqref{eq:6.12}, tabulate monthly condensation-viable
hours ($T_\mathrm{dp} > 10$\,\degC), and aggregate per-module condensate yield. Results
are validated against \cite{alqaed2023awg,eslami2021tec,shourideh2021tec}.

\subsection{Phase 2 --- System Design}

Mechanical concept designs for all three scales; water treatment train sizing (filter
stages, UV dose, pump duty); power system sizing via Eqs.~\eqref{eq:6.15}--\eqref{eq:6.19};
control state machine (Harvest / Standby / Clean / Emergency) to be implemented on
an ESP32-class microcontroller.

\subsection{Phase 3 --- Prototype Preparation (REE 598 Conclusion)}

Scale~B BOM with supplier quotes; fabrication sequence; validation test plan (condensate
yield, TDS, contact force, battery SoC, transmittance restoration, 72-h autonomous cycle).

\section{Computer-Aided Engineering Tools}

\begin{itemize}[leftmargin=*]
  \item \textbf{Python (NumPy, Pandas, Matplotlib)}: psychrometric simulation, TEC
        operating-point solver, power budget calculations.
  \item \textbf{PVsyst}: annual PV yield and monthly soiling sensitivity simulation.
  \item \textbf{MATLAB / Simulink}: control system model and state-machine simulation.
  \item \textbf{AutoCAD / SolidWorks}: mechanical design and dimensioned drawings.
  \item \textbf{LTspice}: DC power system simulation (MPPT, battery, load scheduling).
\end{itemize}

\section{Preliminary Results and Analysis}

\subsection{Condensate Yield}

Table~\ref{tab:prelim_yield} compares modelled yields against literature ranges and the
design targets from Eq.~\eqref{eq:3.2}.

\begin{table}[H]
\centering
\caption{Preliminary condensate yield: thermodynamic model vs.\ literature vs.\ target.}
\label{tab:prelim_yield}
\begin{tabular}{lcccc}
\toprule
\textbf{Scale} & \textbf{Modules} & \textbf{Model (L/day)} &
\textbf{Literature range (L/day)} & \textbf{Target (L/day)} \\
\midrule
A & 1  & 0.155 & 0.07--0.27  & \ScaleAWaterDay \\
R & 6  & 0.93  & 0.42--1.62  & \ScaleRWaterDay \\
B & 12 & 1.86  & 0.84--3.24  & \ScaleBWaterDay \\
C & 80 & 12.4  & 5.6--21.6   & \ScaleCWaterDay \\
\bottomrule
\multicolumn{5}{l}{\footnotesize Summer design point: $T_\mathrm{amb}$ = \TambSummer\,\degC,
RH = \RHdesign\,\%, $t_\mathrm{harvest}$ = 6\,h. At winter conditions (RH = 55\,\%,
$t_\mathrm{harvest}$ = 10\,h),}\\
\multicolumn{5}{l}{\footnotesize per-module yield rises to $\approx$0.42\,L/day, resolving the
Scales~R/B/C gap without additional hardware.}
\end{tabular}
\end{table}

\subsection{Preliminary Economic Analysis}

\begin{table}[H]
\centering
\caption{Preliminary economic analysis by deployment scale (reconciled with Table~\ref{tab:energy_recovery} using the 20\,\% untreated-soiling annual loss assumption).}
\label{tab:economics}
\begin{tabular}{lcccc}
\toprule
\textbf{Parameter}                & \textbf{Scale A} & \textbf{Scale R} & \textbf{Scale B} & \textbf{Scale C} \\
\midrule
Estimated BOM cost (USD)          & 150      & 500      & 900      & 12{,}000 \\
Annual energy recovered (kWh)     & 27       & 1{,}700  & 3{,}360  & 67{,}200 \\
Annual benefit at tariff (USD)    & 0.54     & 34.0     & 67.2     & 1{,}344 \\
Simple payback (years)            & 278$^\star$ & 14.7 & 13.4     & 8.9 \\
\bottomrule
\end{tabular}
\vspace{0.25em}
\par\noindent\footnotesize BOM costs revised from the engineering review (Apr.\ 2026): 1.8~kWh LFP at 450~USD market price, commercial-grade heat sinks, 100~L HDPE tank for Scale~B, decomposition of Scale~C into a parallel fleet. Tariff \TariffUSD\,USD/kWh \cite{egyptera_tariff}. $^\star$Scale~A payback is fleet-viable only --- a single street-lamp unit generates only 27\,kWh/yr of recovery; a New-Administrative-Capital-scale deployment of $\sim$10{,}000 lamps produces 270\,MWh/yr of recovered output and amortises fleet logistics and shared fabrication tooling. \normalsize
\end{table}

\noindent Scales~R, B, and C all show payback periods within the 10--15-year window
typical of PV O\&M hardware. Scale~R (residential) is particularly attractive given
the prevalence of villa rooftops across Egypt's new cities and the low labour cost of
retrofit. Scale~A (street-lamp) is fleet-viable across thousands of municipal units.

% ============================================================
%  CHAPTER 7 — CONCLUSIONS AND FUTURE WORK
% ============================================================
\chapter{Conclusions and Future Work}

\section{Main Outcomes of REE 598}

This proposal has established:

\begin{enumerate}[leftmargin=*, label=(\roman*)]
  \item \textbf{Wet cleaning with low-TDS TEC condensate} (TDS $<$ 5\,mg/L) is the
        highest-efficiency PV panel cleaning method, restoring 90--98\,\% of optical
        transmittance and eliminating the water-spot residue that degrades mineralised
        water cleaning. The efficiency advantage over dry-only systems is 15--28
        percentage points of transmittance restoration per event.
  \item \textbf{TEC-based AWH is thermodynamically feasible} at Egyptian arid design
        conditions, producing $\approx$0.155\,L/module/day in summer and
        $\approx$0.42\,L/module/day in winter, consistent with published experimental
        data from comparable climates \cite{alqaed2023awg,eslami2021tec,shourideh2021tec}.
  \item \textbf{Autonomous, maintenance-free operation} is the central design
        achievement. The robot requires no water supply connection, no cleaning crew,
        and no daily human presence. The TEC condenser has no moving parts beyond a
        fan, the condensate source is unlimited and atmospheric, and the only consumables
        are filter cartridges replaced every 1--3 months. This transforms PV panel
        cleaning from a labour-intensive operational task into a set-and-forget
        infrastructure function.
  \item \textbf{Self-powered from the maintained PV array,} the robot consumes only
        2--10\,\% of daily output --- a negligible parasitic load that is far outweighed
        by the energy recovered through regular wet cleaning (15--20\,\% annual yield
        improvement relative to no cleaning).
  \item \textbf{Economic viability is confirmed} for Scales~R, B, and C (payback
        8.9--14.7 years under the reconciled 20\,\% untreated-soiling model) and is
        fleet-viable for Scale~A across Egyptian new-city street-lamp deployments.
        The priority deployment order is \emph{street-lamp (A) $\to$ residential (R)
        $\to$ commercial (B) $\to$ industrial (C, stretch case)}.
\end{enumerate}

\section{Tasks to Be Accomplished in REE 599 (Senior Design Project II)}

\begin{enumerate}[leftmargin=*, label=\textbf{T\arabic*.}]
  \item Fabricate the Scale~B prototype (BOM in Appendix~\ref{app:bom}).
  \item Conduct a 72-hour outdoor autonomous cycle test on a 10-panel test rig.
  \item Measure and log: condensate yield vs.\ ambient RH; post-treatment TDS; brush
        contact force; battery SoC trajectory; transmittance restoration.
  \item Validate all measurements against model predictions
        (Eqs.~\eqref{eq:6.4}--\eqref{eq:6.13}).
  \item Refine module count for Scales~R, B, and C based on measured yield data.
  \item Produce engineering drawings: Scale~A, R, and B robot assembly, TEC mounting
        detail, water treatment schematic, control wiring diagram.
  \item Complete and submit the REE~599 final report with recommendations for
        commercial deployment in Egyptian new cities.
\end{enumerate}

\section{Project Gantt Chart}

\begin{landscape}
\begin{figure}[H]
\centering
\begin{ganttchart}[
  hgrid,
  vgrid,
  x unit = 0.70cm,
  y unit title = 0.56cm,
  y unit chart = 0.50cm,
  title label font=\small,
  bar label font=\scriptsize,
  group label font=\small\bfseries,
  bar/.append style={fill=blue!40},
  group/.append style={fill=blue!65,draw=blue!65},
  milestone/.append style={fill=red!60,draw=red!60},
  title height=1
]{1}{20}
  \gantttitle{REE 598 (Weeks 1--14) \quad REE 599 (Weeks 15--20)}{20} \\
  \gantttitlelist{1,...,20}{1} \\
  \ganttgroup{REE 598 --- Phase 1: Climate \& TEC Feasibility}{1}{4} \\
  \ganttbar{NASA POWER data retrieval \& processing}{1}{1} \\
  \ganttbar{Psychrometric simulation (Python)}{1}{2} \\
  \ganttbar{TEC thermodynamic model \& validation}{2}{4} \\
  \ganttgroup{REE 598 --- Phase 2: System Design}{4}{10} \\
  \ganttbar{Robot mechanical concept (3 scales)}{4}{6} \\
  \ganttbar{Water treatment train sizing}{5}{6} \\
  \ganttbar{Power system \& battery sizing}{6}{7} \\
  \ganttbar{Control state machine design}{7}{8} \\
  \ganttbar{PVsyst soiling simulation}{6}{9} \\
  \ganttbar{CAD concept drawings (Scales A \& B)}{8}{10} \\
  \ganttmilestone{System design review}{10} \\
  \ganttgroup{REE 598 --- Phase 3: Prototype Preparation}{10}{14} \\
  \ganttbar{Scale B BOM \& supplier quotes}{10}{11} \\
  \ganttbar{Fabrication plan \& test protocol}{11}{12} \\
  \ganttbar{Procurement --- long-lead items}{12}{14} \\
  \ganttmilestone{REE 598 report submission}{14} \\
  \ganttgroup{REE 599 --- Phase 4: Prototype Build \& Validation}{14}{20} \\
  \ganttbar{TEC subassembly \& bench test}{14}{15} \\
  \ganttbar{Robot chassis fabrication}{15}{16} \\
  \ganttbar{Electronics integration \& firmware}{16}{17} \\
  \ganttbar{Field test (72-h autonomous cycle)}{17}{18} \\
  \ganttbar{Data analysis \& final report}{18}{20} \\
  \ganttmilestone{Final prototype validation}{20} \\
\end{ganttchart}
\caption{Project Gantt chart: REE 598 Phases 1--3 and REE 599 Phase 4.}
\label{fig:gantt}
\end{figure}
\end{landscape}

% ============================================================
%  REFERENCES
% ============================================================
\bibliographystyle{ieeetr}
\bibliography{references}

% ============================================================
%  APPENDICES
% ============================================================
\appendix

\chapter{Python Simulation Outputs}
\label{app:python}

The following files are generated by the project computation scripts. Run scripts from
the project root directory before compiling.

\begin{itemize}[leftmargin=*]
  \item \texttt{output/nasa\_power\_hourly.csv} --- five-year hourly NASA POWER data
        for Benban.
  \item \texttt{output/psychrometrics\_monthly.csv} --- monthly dew-point distribution,
        condensation hours, and per-module yield statistics.
  \item \texttt{output/calculation\_summary.txt} --- TEC sizing, power budget, and
        economic results.
  \item \texttt{output/lcow\_table\_autogen.tex} --- AWH technology LCOW/SEC
        comparison table (retained as supporting reference material).
\end{itemize}

\chapter{TEC1-12706 Module Design Parameters}
\label{app:tec}

\begin{table}[H]
\centering
\caption{TEC1-12706 module parameters used in design calculations
(Appendix~\ref{app:tec}).}
\label{tab:A2}
\begin{tabular}{llc}
\toprule
\textbf{Parameter}       & \textbf{Symbol}            & \textbf{Value} \\
\midrule
Maximum current          & $I_\mathrm{max}$            & \TECImax\,A \\
Maximum voltage          & $V_\mathrm{max}$            & \TECVmax\,V \\
Maximum heat pumping     & $Q_\mathrm{max}$ (at $T_h$=27\,\degC)  & \TECQmax\,W \\
Maximum $\Delta T$       & $\Delta T_\mathrm{max}$     & \TECDTmax\,K \\
Seebeck coefficient      & $\alpha$                    & \TECAlpha\,V/K \\
Internal resistance      & $R$                         & \TECR\,$\Omega$ \\
Thermal conductance      & $K$                         & \TECK\,W/K \\
Figure of merit          & $Z$                         & $2.76\times10^{-3}$\,K$^{-1}$ \\
Module dimensions        & ---                         & 40\,$\times$\,40\,$\times$\,3.4\,mm \\
\bottomrule
\end{tabular}
\end{table}

\chapter{Scale B Prototype Bill of Materials}
\label{app:bom}

\begin{table}[H]
\centering
\caption{Scale~B prototype bill of materials (indicative prices, April 2026).}
\label{tab:A3}
\begin{tabular}{p{4.5cm}p{4.0cm}rr}
\toprule
\textbf{Component} & \textbf{Specification} & \textbf{Qty} & \textbf{Unit cost (USD)} \\
\midrule
TEC module                 & TEC1-12706, 40$\times$40\,mm & 12 & 3.50 \\
Heat sink                  & Al extrusion, 100$\times$60\,mm, 1.5\,K/W-grade & 12 & 8.50 \\
Cooling fan                & 60\,mm, 12\,V, 1.0\,W        & 12 & 2.80 \\
Cold plate                 & Al anodised, 40$\times$40$\times$3\,mm & 12 & 1.50 \\
Condensate tank            & HDPE sealed, 100\,L           &  1 & 35.00 \\
Sediment filter            & 5\,\textmu m cartridge        &  1 & 5.00 \\
Activated carbon filter    & Standard 10'' cartridge       &  1 & 7.00 \\
UV-C sterilisation module  & 3\,W, 254\,nm, in-line        &  1 & 18.00 \\
Diaphragm pump             & 12\,V, 2\,bar, 3\,L/min       &  1 & 22.00 \\
Fan spray nozzles          & 1/4'' BSP                     &  8 & 2.50 \\
Flow meter                 & Pulse, 1--10\,L/min           &  1 & 9.00 \\
Microcontroller            & ESP32 DevKit V1               &  1 & 4.50 \\
Temperature/RH sensors     & DHT22                         &  4 & 3.50 \\
Current sensors            & ACS712-20A                    &  2 & 2.20 \\
Level sensor               & Ultrasonic, 12\,V             &  1 & 6.00 \\
DC buck converter          & 24\,V to TEC, 10\,A           &  1 & 12.00 \\
MPPT charge controller     & 20\,A, 24\,V                  &  1 & 35.00 \\
LiFePO$_4$ battery         & 2500\,Wh, 24\,V / 100\,Ah (with BMS) &  1 & 420.00 \\
Robot chassis              & Al extrusion, wheels, motors  &  1 & 95.00 \\
Microfibre brush bar       & 250\,mm wide                  &  1 & 18.00 \\
Wiring, connectors, tubing & Assorted                      & -- & 30.00 \\
\midrule
\multicolumn{3}{l}{\textbf{Total estimated cost}} & \textbf{$\approx$\,USD\,900} \\
\bottomrule
\end{tabular}
\end{table}

\chapter{Relevant Standards and Codes}
\label{app:standards}

\begin{itemize}[leftmargin=*]
  \item \textbf{IEC 61215:2021} --- Terrestrial PV modules: design qualification and
        type approval.
  \item \textbf{IEC 62446-1:2016} --- PV systems: requirements for testing,
        documentation and maintenance.
  \item \textbf{IEC 60529:2013} --- Degrees of protection (IP Code).
  \item \textbf{EgyptERA Tariff Circular, August 2024} --- Electricity tariff for
        industrial and commercial consumers \cite{egyptera_tariff}.
  \item \textbf{Egyptian National Building Code, Part 7} --- Renewable energy system
        installation requirements for new cities.
\end{itemize}

\end{document}
